\section*{Chapitre \ref{ch:s2:systematic-credit-and-bond-etf-arbitrage} --- Arbitrage systématique sur crédit et arbitrage sur ETF obligataires}

\begin{solution}{exo:s2:systematic-credit-and-bond-etf-arbitrage:1}
Environ 300 points de base~: le cours vaut $0.9496$ de la valeur liquidative (VL) et celle-ci vaut $1.0215$ de la valeur vraie, donc le cours vaut $0.9496 \times 1.0215 = 0.970$ de la
valeur vraie. Les deux composantes se multiplient au lieu de s'additionner, ce qui explique pourquoi $215 + 300$ dépasse 504.
\end{solution}

\begin{solution}{exo:s2:systematic-credit-and-bond-etf-arbitrage:2}
Acheter à $0.97$ de la valeur vraie et vendre les obligations à leur valeur vraie rapporte $1/0.97 - 1 = 3.09~\%$ de la valeur rachetée~; après déduction du coût de 1.5~\%, cela fait
159 points de base. Le meilleur jour du chapitre, 163, ajoute le bruit de prix.
\end{solution}

\begin{solution}{exo:s2:systematic-credit-and-bond-etf-arbitrage:3}
$1.50~\% \times 6.5 = 9.75~\%$, proche du $9.8~\%$ simulé~; le carry sur trois semaines et la dispersion des durations autour de 6.5 expliquent le reste.
\end{solution}

\begin{solution}{exo:s2:systematic-credit-and-bond-etf-arbitrage:4}
Le rendement excédentaire d'une obligation dépend surtout de sa duration de spread multipliée par la variation de son spread, et les spreads augmentent avec la duration. Un signal comparé à l'ensemble des obligations
achèterait des obligations longues pour leurs spreads élevés et s'exposerait au risque de crédit de marché, qui est un bêta et non un facteur. Une comparaison avec des pairs de
duration et de notation similaires, et un équilibrage des positions longues et courtes en duration de spread, isolent la caractéristique.
\end{solution}

\begin{solution}{exo:s2:systematic-credit-and-bond-etf-arbitrage:5}
Haddad, Moreira et Muir ont constaté que les investisseurs ont tenté de lever des liquidités en vendant leurs avoirs les plus sûrs et les plus liquides. Les ventes
ont davantage touché les ETF sur bons du Trésor, obligations municipales et obligations investment grade, détenus comme réserves liquides, que les ETF high yield.
\end{solution}

\begin{solution}{exo:s2:systematic-credit-and-bond-etf-arbitrage:6}
L'indice détient des milliers d'obligations, dont beaucoup se négocient rarement~; les livrer ou les recevoir toutes serait coûteux, voire impossible. Koont, Ma, Pastor
et Zeng ont montré que les ETF obligataires composent leurs paniers avec du cash et un sous-ensemble d'obligations, d'autant plus que celles-ci sont peu liquides, et les ajustent pour corriger les déséquilibres
entre créations et rachats tout en permettant l'arbitrage.
\end{solution}

\begin{solution}{exo:s2:systematic-credit-and-bond-etf-arbitrage:7}
Lorsque les obligations ne se négocient que 5~\% des jours, la pire décote double pour atteindre 1\,050 points de base, encore quatorze jours de négociation après le début du stress, et la part périmée
passe de 215 à 838 points de base~: la VL retarde sur la valeur vraie d'environ dix-neuf jours en moyenne, plus longtemps que le stress, si bien qu'elle manque l'essentiel de
l'élargissement. La part due à la pression vendeuse ne change pas, puisque le cours suit la valeur vraie, mais une décote par rapport à la VL mesure désormais surtout des cours
périmés.
\end{solution}

\begin{solution}{exo:s2:systematic-credit-and-bond-etf-arbitrage:8}
La VL est calculée à partir de cours périmés~; dans le stress du chapitre, 215 des 504 points de base correspondent à une VL qui n'a pas rattrapé la baisse des obligations, et
l'écart se résorbe par la baisse de la VL et non par la hausse du cours. Seule la part par rapport à la valeur vraie, 300 points de base, constitue un rendement, et seulement pour
celui qui peut racheter et vendre les obligations, à un coût de 150 points de base pendant le stress.
\end{solution}

\begin{solution}{pb:s2:systematic-credit-and-bond-etf-arbitrage:1}
\begin{enumerate}
\item Une caractéristique des obligations, comme le spread par rapport aux pairs, le momentum de spread, la duration ou la taille, qui les trie en portefeuilles dont les rendements excédentaires ajustés de la duration sont durablement
différents.
\item Les portefeuilles taille, faible risque, valeur et momentum ont dégagé des alphas significatifs tant sur le plan économique que statistique~; de faibles corrélations ont rendu une
combinaison plus efficace~; les résultats tiennent après coûts et sur les obligations liquides.
\item 300 obligations sur huit ans, des durations allant de un à douze~; les spreads évoluent avec un facteur de marché, une tendance persistante par émetteur, un écart qui revient vers sa moyenne et
du bruit~; les rendements excédentaires attendus croissent avec la racine carrée de la duration. La valeur est le spread au-dessus de la droite ajustée de spread en fonction de la duration, moins
sa variation sur trois mois~; le momentum est le resserrement sur six mois~; le faible risque achète des obligations courtes contre des obligations longues à risque égal.
\item Valeur 0.46, momentum 0.14, faible risque 0.95, combiné 0.94~; corrélations de $-0.09$ (valeur, momentum), $-0.03$ (valeur, faible risque) et $-0.02$
(momentum, faible risque).
\item Arbitrage création--rachat~: livrer un panier contre de nouvelles parts lorsque les parts se négocient cher, ou les parts contre le panier lorsqu'elles se négocient bon marché, en gardant l'écart
diminué des coûts. La décote~: cours inférieur à la VL publiée.
\item Les ETF obligataires composent leurs paniers avec du cash et un sous-ensemble de l'indice, surtout lorsque les obligations sont peu liquides, et les ajustent pour corriger les déséquilibres~; l'inclusion
dans le panier favorise la liquidité en général, mais la dégrade lors de déséquilibres importants, comme la crise de la COVID-19.
\item La moyenne des dernières valeurs négociées des obligations~; chaque obligation se négocie un jour sur quatre.
\item Une transaction portant sur une liste entière d'obligations à un prix unique négocié, évaluée par rapport à l'ETF et aux cotations des obligations, et couvrable avec l'ETF.
\item Six ans après le début, les spreads de marché s'élargissent de 150 points de base sur quinze jours, la moitié se résorbe sur quarante jours~; le cours de l'ETF baisse encore de 3~\% sous
la valeur vraie au plus profond et se redresse sur quarante jours. Le panier perd 9.8~\%.
\item Au pire, $-504$ points de base par rapport à la VL~: VL périmée de 215 au-dessus de la valeur vraie et cours 300 en dessous, multipliés entre eux.
\item Sur 28 jours, tous dans le stress, un gain allant jusqu'à 163 points de base de la valeur rachetée après un coût de 1.5~\%~; jamais en dehors du stress, où la
décote moyenne est de 2.7 points de base, avec un écart type de 21.6.
\item Les ETF obligataires liquides se sont négociés avec de fortes décotes par rapport à la VL en mars 2020, davantage pour les bons du Trésor, les obligations municipales et investment grade que pour le high yield, car les investisseurs vendaient leurs actifs les plus sûrs pour obtenir du cash~; les
perturbations se sont inversées après les annonces de la Fed des 23 mars et 9 avril.
\item Ratios de Sharpe de 0.95 (faible risque), 0.46 (valeur), 0.14 (momentum) et 0.94 combiné~; dans le stress, l'ETF s'est négocié 5.0~\% sous la VL et 3.0~\% sous
la valeur vraie, et le participant autorisé pouvait gagner sur 28 jours jusqu'à 1.63~\%.
\item La valeur vraie est la bonne, mais personne ne la voit, car la plupart des obligations n'ont pas été négociées~: la VL est périmée et le cours est poussé par les vendeurs.
\item L'écart qui revient vers sa moyenne explique une part de chaque variation de spread sur six mois, si bien que le momentum achète en partie des obligations dont l'écart est sur le point de revenir, à l'encontre de la valeur.
\item À partir du cours de l'ETF, des cours des obligations qui se sont négociées et de leur co-mouvement avec celles qui ne se sont pas négociées, ainsi que des cotations des dealers, avec un modèle de l'ampleur
du mouvement de la pression vendeuse sur l'ETF.
\item Par rapport à une estimation de la valeur vraie à partir des cours négociés, et non de la VL périmée~; avec les règles de panier telles qu'elles étaient et les coûts de vente des obligations en période de stress.
\item L'arbitrage prime--décote sur ETF.
\item La base négative du chapitre~18 correspond à des obligations en dessous de leur CDS, parce que le financement est rare~; la décote de ce chapitre correspond à un ETF en dessous de ses obligations, parce que le cash est rare~; les
deux s'élargissent dans les mêmes crises et profitent à celui qui dispose encore de bilan.
\item Le chemin parcouru par la valeur vraie des obligations depuis leur dernière négociation, auquel s'ajoute ce que les vendeurs de l'ETF acceptent de payer pour l'immédiateté.
\end{enumerate}
\end{solution}

\begin{solution}{iq:s2:systematic-credit-and-bond-etf-arbitrage:1}
Lorsque les parts se négocient au-dessus de la valeur du panier, un participant autorisé achète le panier, le livre contre de nouvelles parts et les vend~; en dessous, le
participant achète des parts, les rachète contre le panier et vend ce dernier. Les deux transactions poussent le cours vers la valeur du panier, à la
limite des coûts du participant.
\end{solution}

\begin{solution}{iq:s2:systematic-credit-and-bond-etf-arbitrage:2}
Le spread par rapport à une courbe ajustée sur des pairs de notation, de duration et de secteur similaires, éventuellement par rapport à un modèle de risque de défaut~; retirer la part qui correspond à un
élargissement récent, équilibrer les positions en duration de spread, vérifier qu'il ne s'agit pas seulement de risque de défaut, et valoriser le tout avec des coûts réalistes sur les obligations.
\end{solution}

\begin{solution}{iq:s2:systematic-credit-and-bond-etf-arbitrage:3}
Par rapport à l'ETF et aux cotations des obligations~: couvrir la liste immédiatement avec l'ETF, estimer le temps nécessaire pour dénouer chaque obligation et son coût, ajouter
le risque que la liste diverge de l'ETF pendant ce temps, et demander une prime plus élevée pour les obligations que l'ETF ne détient pas.
\end{solution}

\begin{solution}{iq:s2:systematic-credit-and-bond-etf-arbitrage:4}
L'ETF passe sous ses obligations et le livre est valorisé avec une perte~; les obligations ne peuvent pas être vendues à leurs valeurs de marquage, les rachats livrent des paniers qui peuvent ne pas
correspondre au livre, et le financement se resserre. La couverture ne fonctionne que pour celui qui peut tenir jusqu'à ce que la décote se résorbe.
\end{solution}

\begin{solution}{iq:s2:systematic-credit-and-bond-etf-arbitrage:5}
Partir de la dernière transaction de chaque obligation et des cotations des dealers~; la déplacer selon les mouvements observés de proxies liquides (l'ETF, les CDS, les obligations du même émetteur
et des obligations similaires qui se sont négociées), pondérées par sa duration et sa notation~; signaler son ancienneté et son niveau de confiance, et tester l'estimation contre la transaction suivante.
\end{solution}

\begin{solution}{iq:s2:systematic-credit-and-bond-etf-arbitrage:6}
La dernière transaction remonte à $k$ jours avec la probabilité $p(1-p)^k$, si bien que la VL reflète une valeur vraie vieille de $k$ jours, et $E[k] = \sum_k k\, p(1-p)^k =
(1-p)/p$. Après un mouvement régulier de $m$ par jour, qui a duré plus longtemps que les écarts entre transactions de la plupart des obligations, l'écart périmé vaut $m\,E[k] = m(1-p)/p$~:
trois jours de mouvement pour $p = 1/4$, dix-neuf pour $p = 1/20$.
\end{solution}
